% =============================================================
%  Chapter 4 — Material-Aware Risk-Sensitive
%              Continual Hamiltonian Learning
% =============================================================

\chapter{Material-Aware Risk-Sensitive Continual Hamiltonian Learning}
\label{chap:material}

\providecommand{\Ham}{\mathcal{H}}
\providecommand{\Potential}{\mathcal{R}}
\providecommand{\PhaseState}{z}
\providecommand{\Pos}{q}
\providecommand{\Mom}{p}
\providecommand{\Mass}{\mathrm{M}}
\providecommand{\J}{\mathbb{J}}

% -----------------------------------------------------------------------
\section{Why a Richer Learner Is Needed}
\label{sec:ch4:motivation}
% -----------------------------------------------------------------------

Chapter~\ref{chap:grlsnam} established that a geometry-only Hamiltonian
learner can navigate effectively under local sensing: it avoids obstacles,
reaches goals efficiently, and adapts its force-field coefficients online.
Yet three structural gaps remain, each arising from the same root cause:
the potential $\Potential(\vec{q};\mathcal{E})$ responds \emph{only} to
signed geometric clearances.

A region of soft mud from geometrically traversable but materially dangerous terrain 
that immobilizes the robot, a shallow river that disables electronics,
a railway crossing invisible to LiDAR receives the same energy as clear
grass.
The force field has no channel for material risk, no tail-risk objective,
and no coefficient slot for semantic enrichment.
These are not tuning failures; they are architectural gaps.

The challenge of this chapter is therefore:

\begin{quote}\itshape
Not to replace the geometry-only framework, but to \textbf{enlarge it}
along four coupled axes of stored energy, coefficient space, objective, and
adaptation hierarchy so that the same structured phase-space learner also
responds to latent risk structure while remaining stable, smooth, and
progressively valid over a broader configuration space.
\end{quote}

Each axis of enlargement is not independent: adding a new energy term
$H_{\mathrm{mat}}$ induces a new force channel $F_{\mathrm{mat}}$, which
requires a new coefficient $(\lambda_s, \lambda_h)$, which opens a new
gradient pathway in the learning law, which demands a new timescale in
the adaptation hierarchy.
The rest of this chapter gives the mathematics of that chain, from the
enriched Hamiltonian through the gradient-field learning law to the
three-loop curriculum optimizer.

% -----------------------------------------------------------------------
\section{Enriched Problem Setting}
\label{sec:ch4:setting}
% -----------------------------------------------------------------------

\paragraph{Same task, richer context.}
The navigation task is identical to Chapter~\ref{chap:grlsnam}:
reach $\vec{x}_g$ from $\vec{q}_0$ through an unknown workspace.
The enrichment lies entirely in the environment context
$\mathcal{E}_\tau$.
In the geometry-only case, $\mathcal{E}_\tau$ contained only obstacle
clearances.
Here it additionally contains:

\begin{equation}
  \mathcal{E}_\tau
  = \bigl\{
      \vec{c}_{\mathrm{target}},\;
      \{d_i\}_{i=1}^{C},\;
      \underbrace{\tilde{r}(\vec{c})}_{\text{soft risk field}},\;
      \underbrace{\phi(\vec{c})}_{\text{hard hazard SDF}},\;
      \underbrace{P_t(\vec{q}) \in \mathbb{R}^{6\times P\times P}}_{\text{local risk patch}}
    \bigr\},
  \label{eq:ch4:context}
\end{equation}
where $\tilde{r} : \mathcal{Q}_f \to [0,1]$ is a continuous soft-risk
field (e.g., material-class risk score), and
$\phi : \mathcal{Q}_f \to \mathbb{R}$ is the signed-distance function
to hard hazard zones $\mathcal{H}$ (regions where contact is catastrophic).
The local risk patch $P_t(\vec{q})$ is a $6$-channel crop
re-centred at $\vec{q}$ on every integration step, carrying
$(\tilde{r}, \partial_x\tilde{r}, \partial_y\tilde{r},
   h_{\mathrm{mask}}, \phi, \|\nabla\phi\|)$.

\paragraph{Oracle material access.}
In the current chapter, $\tilde{r}$ and $\phi$ are provided by an oracle
(pre-computed from ground-truth material maps, e.g., the DFC2018 Houston
aerial dataset).
This matches \emph{Setting 2} of the thesis scope: the learner has true
material information and must learn to use it; perception-led material
inference from raw imagery is the natural next step, deferred to future
work.

Fig.~\ref{fig:teaser} summarizes our setting and motivation: under partial observability, a static prior is hazard-blind, while our belief-guided structured learner uses local RGB--IR observations to infer material risk and improve navigation over repeated experience.

\begin{figure*}[htbp]
    \centering
    \includegraphics[width=0.75\textwidth]{figs/teaser.png}
    \caption{
    \textbf{From hazard-blind planning to belief-guided structured navigation.}
    Static planners take short but unsafe routes because material hazards are latent. 
    Our continual learner uses local RGB--IR evidence to infer uncertain material beliefs, 
    combines them with geometry and dissipation in a factored structured flow, and improves 
    with experience into safer, smoother navigation.
    }
    \label{fig:teaser}
\end{figure*}

% -----------------------------------------------------------------------
\section{Hierarchy of Complexity: Four Coupled Enrichments}
\label{sec:ch4:hierarchy}
% -----------------------------------------------------------------------

The core thesis claim is that the geometry-only learner is enlarged in
\emph{four coupled ways}, not in one.
Table~\ref{tab:ch4:hierarchy} makes this explicit.

\begin{table}[t]
\centering
\caption{%
  \textbf{Hierarchy of enrichments from Stage~1 to Stage~2.}
  Each new energy term induces a chain of consequences through
  force channels, coefficients, objective terms, and update laws.
  This chain is not cosmetic: removing any link breaks the others.
}
\label{tab:ch4:hierarchy}
\small
\setlength{\tabcolsep}{5pt}
\renewcommand{\arraystretch}{1.35}
\resizebox{\textwidth}{!}{%
\begin{tabular}{lllll}
\toprule
\textbf{Added term} &
\textbf{New force channel} &
\textbf{New coefficient} &
\textbf{New objective term} &
\textbf{New update law} \\
\midrule
$H_{\mathrm{geom}}$ (Ch.~\ref{chap:grlsnam})
  & $F_{\mathrm{goal}}, F_{\mathrm{bar}}$
  & $\beta, \{\alpha_i\}$
  & $\mathcal{L}_{\mathrm{traj}}, \mathcal{L}_{\mathrm{clear}}$
  & Online $\alpha_i$ update \\
\midrule
$H_{\mathrm{mat}}^{\mathrm{soft}}$
  & $F_{\mathrm{soft}} = -\lambda_s\nabla\tilde{r}$
  & $\lambda_s$
  & $w_r\sum_t \tilde{r}(\vec{c}_t)\Delta s_t$
  & Episode CVaR grad.\  \\
$H_{\mathrm{mat}}^{\mathrm{hard}}$
  & $F_{\mathrm{hard}} = -\lambda_h\nabla b(\phi)$
  & $\lambda_h$
  & $w_h\sum_t \mathbf{1}[\phi(\vec{c}_t)<1\,\text{m}]$
  & Episode CVaR grad.\  \\
$R_\theta$ (dissipative port)
  & $-R_\theta = -\gamma\Mass^{-1}\Mom$
  & $\gamma$
  & $\mathcal{L}_{\mathrm{fric}} = |\gamma-\gamma_0|^2$
  & Friction alignment \\
\midrule
\textbf{Total (Stage 2)}
  & $F_{\mathrm{tot}} = F_{\mathrm{goal}}+F_{\mathrm{bar}}$
  & $(\beta,\lambda_s,\lambda_h,$
  & $\mathcal{L}_{\mathrm{P2}} = \mathcal{L}_{\mathrm{nav}}$
  & 3-timescale loop \\
  & $\quad + F_{\mathrm{soft}}+F_{\mathrm{hard}}-R_\theta$
  & $\quad\{\alpha_i\},\gamma)$
  & $\quad+0.3\,\mathcal{L}_{\mathrm{geom}}+w_\lambda\mathcal{L}_\lambda$
  & (§\ref{sec:ch4:multiscale}) \\
\bottomrule
\end{tabular}}
\end{table}

The key message of Table~\ref{tab:ch4:hierarchy} is that this is
\emph{not a simple extension}: adding $H_{\mathrm{mat}}^{\mathrm{soft}}$
is not just "another loss term."
It creates a new force acting on every integration step, a new coefficient
that must be learned, a new gradient pathway through the rollout, and a new
term in the objective that requires a tail-risk criterion rather than an
expected-value surrogate.

% -----------------------------------------------------------------------
\section{Port-Hamiltonian State and Observer Processes}
\label{sec:ch4:portham}
% -----------------------------------------------------------------------

\subsection{Phase-Space Dynamics with Explicit Dissipation}
\label{subsec:ch4:dynamics}

The phase-space state is $\PhaseState_t = (\vec{q}_t, \Mom_t)$ as before.
The key departure from Chapter~\ref{chap:grlsnam} is that dissipation
is now a \emph{separate port}, not buried inside the Hamiltonian.
The discrete-time dynamics are:
\begin{align}
  \Mom_{t+1}
  &= \Mom_t
     - \tau\,\nabla_{\vec{q}} H_\theta(\vec{q}_t,\Mom_t)
     - \tau\,R_\theta(\vec{q}_t,\Mom_t),
  \label{eq:ch4:p_update} \\
  \vec{q}_{t+1}
  &= \vec{q}_t + \tau\,\Mass^{-1}\Mom_{t+1},
  \label{eq:ch4:q_update}
\end{align}
where the \textbf{dissipative port} is:
\begin{equation}
  R_\theta(\vec{q},\Mom) = \gamma\,\Mass^{-1}\Mom,
  \qquad \gamma \geq 0.
  \label{eq:ch4:dissipation}
\end{equation}
This separation is important.
In a port-Hamiltonian system, the stored energy $H_\theta$ and the
dissipation $R_\theta$ play distinct roles: $H_\theta$ defines the
conservative flow topology (basins, barriers), while $R_\theta$ removes
excess kinetic energy.
Conflating them would make the stability analysis unclear and the
gradient pathways entangled.

The energy balance along any trajectory is:
\begin{equation}
  \Delta H_t
  := H_\theta(\vec{q}_{t+1},\Mom_{t+1}) - H_\theta(\vec{q}_t,\Mom_t)
  = -\tau(\Mom_t)^\top R_\theta(\vec{q}_t,\Mom_t)
    + \tau(\Mom_t)^\top \mathbf{B} u_t
  \;\leq\; \tau(\Mom_t)^\top \mathbf{B} u_t,
  \label{eq:ch4:energy_balance}
\end{equation}
which is the discrete passivity condition:
when no external control signal $u_t$ is applied, the stored energy
is non-increasing.
This is monitored as a curriculum criterion (§\ref{sec:ch4:multiscale}).

\subsection{Observer and Progress Monitoring}
\label{subsec:ch4:observer}

The learner actively senses both the environment and its own progress.
The context update between stages is:
\begin{equation}
  y_{t+1} = \Psi(y_t, \vec{q}_{t+1};\, Z^{(1)}, P_{t+1}),
  \label{eq:ch4:observer}
\end{equation}
where $\Psi$ re-centres the local obstacle extractor and risk patch at
$\vec{q}_{t+1}$.
In the oracle setting, $\Psi$ reduces to extracting the new $\{d_i\}$
from $Z^{(1)}$ and reading $\tilde{r}, \phi$ from $Z^{(2)}$ at
$\vec{q}_{t+1}$.

Additionally, the learner monitors two progress scalars at each step:
\begin{equation}
  s_t^{\mathrm{seg}} = |\vec{q}_t - \vec{g}_k|,
  \qquad
  s_t^{\mathrm{pass}} = \Delta H_t,
  \label{eq:ch4:progress}
\end{equation}
where $\vec{g}_k$ is the current segment goal.
$s_t^{\mathrm{seg}}$ tracks spatial progress; $s_t^{\mathrm{pass}}$
tracks energy dissipation (a proxy for dynamical stability).
Both are used in the adaptation loops of §\ref{sec:ch4:multiscale}.

% -----------------------------------------------------------------------
\section{Factored Stored Energy and Force Decomposition}
\label{sec:ch4:energy}
% -----------------------------------------------------------------------

\subsection{Factored Hamiltonian}
\label{subsec:ch4:hamiltonian}

The stored energy decomposes additively:
\begin{equation}
  H_\theta(\vec{q},\Mom)
  = \underbrace{\tfrac{1}{2}\Mom^\top\Mass^{-1}\Mom}_{H_{\mathrm{kin}}}
  + \underbrace{H_{\mathrm{geom}}(\vec{q};\mathcal{E})}_{H_{\mathrm{geom}}}
  + \underbrace{H_{\mathrm{mat}}^{\mathrm{soft}}(\vec{q}) + H_{\mathrm{mat}}^{\mathrm{hard}}(\vec{q})}_{H_{\mathrm{mat}}},
  \label{eq:ch4:ham_factored}
\end{equation}
with $R_\theta$ a \emph{separate} dissipative port (not part of
$H_\theta$).

\paragraph{Geometry sub-Hamiltonian.}
\begin{equation}
  H_{\mathrm{geom}}(\vec{q};\mathcal{E})
  = \beta\,D_{\mathcal{G}}(\vec{q},\vec{g})
    + \sum_{j}\alpha_j\,b_{\mathrm{IPC}}\!\bigl(d_j(\vec{q};\mathcal{E})\bigr),
  \label{eq:ch4:H_geom}
\end{equation}
where $D_{\mathcal{G}}(\vec{q},\vec{g}) = \|\vec{c}-\vec{g}\|^2$ is
goal attraction and $b_{\mathrm{IPC}}$ is an IPC-style log-barrier.
This is the geometry-only potential of Chapter~\ref{chap:grlsnam},
carried forward unchanged.

\paragraph{Soft material-risk sub-Hamiltonian.}
\begin{equation}
  H_{\mathrm{mat}}^{\mathrm{soft}}(\vec{q})
  = \lambda_s\,\tilde{r}(\vec{c}),
  \label{eq:ch4:H_mat_soft}
\end{equation}
where $\tilde{r}(\vec{c}) \in [0,1]$ is the soft-risk field at the
robot's center-of-mass position.
The induced force is $F_{\mathrm{soft}} = -\lambda_s\nabla_{\vec{c}}\tilde{r}$,
which deflects the trajectory away from high-risk terrain.

\paragraph{Hard hazard-barrier sub-Hamiltonian.}
\begin{equation}
  H_{\mathrm{mat}}^{\mathrm{hard}}(\vec{q})
  = \lambda_h\,b_{\mathrm{sp}}\!\bigl(\phi(\vec{c})\bigr),
  \label{eq:ch4:H_mat_hard}
\end{equation}
where $\phi$ is the signed-distance to hard hazard zones $\mathcal{H}$,
and $b_{\mathrm{sp}}(s) = \mathrm{softplus}(-s/\sigma_\phi)$ grows large
as $\phi \to 0$.
The induced force $F_{\mathrm{hard}} = -\lambda_h\nabla_{\vec{c}}b_{\mathrm{sp}}(\phi)$
provides forward invariance of $\{\phi > 0\}$ in the spirit of
Definition~\ref{def:cbf}.

\subsection{Total Force Decomposition}
\label{subsec:ch4:forces}

The total force on the momentum is:
\begin{equation}
  F_{\mathrm{tot}}(\vec{q}_t,\Mom_t)
  = \underbrace{F_{\mathrm{goal}}}_{{-\beta\nabla D_\mathcal{G}}}
  + \underbrace{F_{\mathrm{bar}}}_{{-\sum_j\alpha_j\nabla b_{\mathrm{IPC}}}}
  + \underbrace{F_{\mathrm{soft}}}_{{-\lambda_s\nabla\tilde{r}}}
  + \underbrace{F_{\mathrm{hard}}}_{{-\lambda_h\nabla b_{\mathrm{sp}}}}
  - \underbrace{R_\theta}_{\gamma\Mass^{-1}\Mom},
  \label{eq:ch4:F_tot}
\end{equation}
so the momentum update \eqref{eq:ch4:p_update} becomes
$\Mom_{t+1} = \Mom_t + \tau\,F_{\mathrm{tot}}(\vec{q}_t,\Mom_t)$.

% -----------------------------------------------------------------------
\section{Multi-Term Phase-Space Objective}
\label{sec:ch4:objective}
% -----------------------------------------------------------------------

\subsection{Episode Cost and CVaR Objective}
\label{subsec:ch4:cvar}

The episode cost accumulates over all integration steps:
\begin{equation}
  J(\theta)
  = w_g\,D_{\mathcal{G}}(\vec{q}_T,\vec{g})
  + w_\ell\,\ell
  + w_r\sum_t \tilde{r}(\vec{c}_t)\Delta s_t
  + w_h\sum_t \mathbf{1}[\phi(\vec{c}_t) < 1\,\text{m}],
  \label{eq:ch4:J}
\end{equation}
where $\ell = \sum_t \|\vec{q}_{t+1} - \vec{q}_t\|$ is path length and
$\Delta s_t$ is the step displacement.
The four terms correspond to: terminal goal error, travel efficiency,
cumulative soft-risk exposure, and hard-hazard incursion count.

The primary navigation loss is the \textbf{CVaR at level $\beta=0.95$}:
\begin{equation}
  \mathcal{L}_{\mathrm{nav}}(\theta)
  = \mathrm{CVaR}_{0.95}\!\bigl(J(\theta)\bigr)
  = \mathbb{E}\!\left[J \;\big|\; J \geq \mathrm{VaR}_{0.95}(J)\right].
  \label{eq:ch4:Lnav}
\end{equation}
This optimizes the \emph{worst 5\% of rollouts}: it directly penalizes
rare catastrophic material encounters, not just their expected frequency.
An expected-cost surrogate ($\mathbb{E}[J]$) would underweight these
rare events; CVaR does not.

\subsection{Full Phase Objectives}
\label{subsec:ch4:phase_obj}

The geometry component is retained from Chapter~\ref{chap:grlsnam}:
\begin{equation}
  \mathcal{L}_{\mathrm{geom}}
  = w_t\,\mathcal{L}_{\mathrm{traj}}
  + w_v\,\mathcal{L}_{\mathrm{vel}}
  + w_\gamma\,|\gamma - \gamma_0|^2
  + w_c\,\mathcal{L}_{\mathrm{clear}}
  + w_m\,\mathcal{L}_{\mathrm{multi}}.
  \label{eq:ch4:L_geom}
\end{equation}

The two training phases use:
\begin{align}
  \mathcal{L}_{\mathrm{P1}}
  &= \mathcal{L}_{\mathrm{geom}},
  \label{eq:ch4:LP1} \\
  \mathcal{L}_{\mathrm{P2}}
  &= \mathcal{L}_{\mathrm{nav}}
    + 0.3\,\mathcal{L}_{\mathrm{geom}}
    + w_\lambda\,\mathcal{L}_\lambda,
  \label{eq:ch4:LP2}
\end{align}
where $\mathcal{L}_\lambda = \|\lambda_s\|^2 + \|\lambda_h\|^2$ is a
coefficient regularizer preventing material heads from dominating the
geometry terms prematurely.
The $0.3$ weight on $\mathcal{L}_{\mathrm{geom}}$ in Phase~2 retains
the geometric structure learned in Phase~1 as a scaffold.

% -----------------------------------------------------------------------
\section{Gradient-Field Learning Law}
\label{sec:ch4:gradients}
% -----------------------------------------------------------------------

\begin{quote}\itshape
The policy is the integrated vector field induced by the factored
Hamiltonian, so learning proceeds by reshaping that field rather than
only attaching losses to final trajectories.
\end{quote}

\subsection{Rollout Sensitivity Recursion}
\label{subsec:ch4:sensitivity}

Because the policy is a numerical integrator applied to
$\dot\PhaseState = F_{\mathrm{tot}}(\PhaseState;\theta)$,
the gradient of any rollout-dependent loss requires propagating
sensitivity through the integrator.
Differentiating \eqref{eq:ch4:p_update}--\eqref{eq:ch4:q_update}
with respect to $\theta$:

\begin{align}
  \frac{\partial \Mom_{t+1}}{\partial\theta}
  &= \frac{\partial \Mom_t}{\partial\theta}
  - \tau\!\left(
      \frac{\partial \nabla_{\vec{q}} H_\theta}{\partial \PhaseState_t}
      \frac{\partial \PhaseState_t}{\partial\theta}
      + \frac{\partial \nabla_{\vec{q}} H_\theta}{\partial\theta}
      + \frac{\partial R_\theta}{\partial \PhaseState_t}
        \frac{\partial \PhaseState_t}{\partial\theta}
      + \frac{\partial R_\theta}{\partial\theta}
    \right),
  \label{eq:ch4:dp_dtheta} \\
  \frac{\partial \vec{q}_{t+1}}{\partial\theta}
  &= \frac{\partial \vec{q}_t}{\partial\theta}
  + \tau\,\Mass^{-1}\frac{\partial \Mom_{t+1}}{\partial\theta}.
  \label{eq:ch4:dq_dtheta}
\end{align}

This is a forward-mode sensitivity ODE: initialized at
$(\partial\PhaseState_0/\partial\theta) = 0$,
it propagates how small changes in $\theta$ alter the entire trajectory.
Crucially, because $R_\theta$ is a separate port, its gradient
$\partial R_\theta/\partial\theta$ enters the recursion independently
of $\partial\nabla_{\vec{q}} H_\theta/\partial\theta$.
The two gradient streams are decoupled, which is why separating the
port is architecturally important.

\subsection{Component-wise Gradient Channels}
\label{subsec:ch4:grad_channels}

The five force channels in \eqref{eq:ch4:F_tot} each contribute a
distinct gradient pathway:

\begin{align}
  \frac{\partial\mathcal{L}}{\partial\beta}
  &= \sum_t \frac{\partial\mathcal{L}}{\partial F_{\mathrm{goal},t}}
     \cdot (-\nabla_{\vec{c}} D_\mathcal{G}),
  \label{eq:ch4:dL_dbeta} \\
  \frac{\partial\mathcal{L}}{\partial\alpha_j}
  &= \sum_t \frac{\partial\mathcal{L}}{\partial F_{\mathrm{bar},t}}
     \cdot (-\nabla_{\vec{q}} b_{\mathrm{IPC}}(d_j)),
  \label{eq:ch4:dL_dalpha} \\
  \frac{\partial\mathcal{L}}{\partial\lambda_s}
  &= \sum_t \frac{\partial\mathcal{L}}{\partial F_{\mathrm{soft},t}}
     \cdot (-\nabla_{\vec{c}}\tilde{r}),
  \label{eq:ch4:dL_dlambda_s} \\
  \frac{\partial\mathcal{L}}{\partial\lambda_h}
  &= \sum_t \frac{\partial\mathcal{L}}{\partial F_{\mathrm{hard},t}}
     \cdot (-\nabla_{\vec{c}} b_{\mathrm{sp}}(\phi)),
  \label{eq:ch4:dL_dlambda_h} \\
  \frac{\partial\mathcal{L}}{\partial\gamma}
  &= \sum_t \frac{\partial\mathcal{L}}{\partial(-R_{\theta,t})}
     \cdot (-\Mass^{-1}\Mom_t).
  \label{eq:ch4:dL_dgamma}
\end{align}

Each equation says: the gradient with respect to a coefficient equals the
sum over time of the loss sensitivity to that force channel times the
channel's analytical Jacobian.
Because the Jacobians are explicit (they are $\nabla\tilde{r}$,
$\nabla b_{\mathrm{IPC}}$, etc.), the backpropagation through the
integrator carries interpretable structure rather than opaque chain rules.

% -----------------------------------------------------------------------
\section{Double Regularization and Smoothness of Updates}
\label{sec:ch4:regularization}
% -----------------------------------------------------------------------

Learning to reshape a force field requires not only a task objective but
also regularization of \emph{the field itself}.
Two levels of regularization are used.

\paragraph{Energy-level regularization $\mathcal{R}_H$.}
Penalizes drift in the learned dissipation coefficient and in the
material force magnitudes:
\begin{equation}
  \mathcal{R}_H(\theta)
  = w_\gamma\,|\gamma - \gamma_0|^2
  + w_\lambda\,(\|\lambda_s\|^2 + \|\lambda_h\|^2).
  \label{eq:ch4:R_H}
\end{equation}
The first term anchors the dissipation near a reference value,
preventing under-damping (passivity violation).
The second term prevents material force heads from overpowering the
geometry backbone before the risk field has been sufficiently explored.

\paragraph{Gradient-field regularization $\mathcal{R}_{\nabla H}$.}
Penalizes non-smooth and near-constraint behavior of the force field:
\begin{equation}
  \mathcal{R}_{\nabla H}(\theta)
  = w_c\,\mathcal{L}_{\mathrm{clear}}
  + w_m\,\mathcal{L}_{\mathrm{multi}},
  \label{eq:ch4:R_nablaH}
\end{equation}
where $\mathcal{L}_{\mathrm{clear}}$ penalizes barrier violations along
the rollout and $\mathcal{L}_{\mathrm{multi}}$ trains near-contact
robustness via perturbed short rollouts.

In principle, the smoothest regularizer is the strong-form Jacobian
penalty on the total force field:
\begin{equation}
  \mathcal{R}_{\mathrm{Jac}}(\theta)
  = w_J\,\frac{1}{T}\sum_t
    \|\nabla_{\vec{q}} F_{\mathrm{tot}}(\vec{q}_t,\Mom_t)\|_F^2,
  \label{eq:ch4:R_Jac}
\end{equation}
which penalizes large spatial gradients in the force field itself,
not just in its output.
Computing \eqref{eq:ch4:R_Jac} exactly requires second-order
differentiation through the integrator, which is expensive.
The thesis uses $\mathcal{R}_{\nabla H}$ as a tractable surrogate,
with $\mathcal{R}_{\mathrm{Jac}}$ identified as a direction for future
work.

\paragraph{Total loss.}
\begin{equation}
  \mathcal{L}_{\mathrm{total}}(\theta)
  = \mathcal{L}_{\mathrm{task}}
  + \mathcal{R}_H(\theta)
  + \mathcal{R}_{\nabla H}(\theta),
  \label{eq:ch4:L_total}
\end{equation}
where $\mathcal{L}_{\mathrm{task}} \in \{\mathcal{L}_{\mathrm{P1}},
\mathcal{L}_{\mathrm{P2}}\}$ depending on the training phase.

% -----------------------------------------------------------------------
\section{Multi-Timescale Optimization}
\label{sec:ch4:multiscale}
% -----------------------------------------------------------------------

The material-aware learner operates at three nested timescales.
Figure~\ref{fig:prelim:timescales} from Chapter~\ref{chap:background}
already illustrated their nesting; here we give the concrete update laws.

\subsection{Segment Timescale: Local Coefficient Alignment}
\label{subsec:ch4:segment}

Within each navigation segment $k$, the response signal
$\Delta \vec{o}_k^{\mathrm{des}} = \vec{o}_k^{\mathrm{des}} - \vec{o}_k$
(desired vs. observed state at the segment endpoint) drives a local
Gauss-Newton update on the current coefficient vector $\theta_k$:
\begin{equation}
  \Delta\theta_k
  = (J_k^\top J_k + \lambda_J I)^{-1} J_k^\top\,\Delta\vec{o}_k^{\mathrm{des}},
  \qquad
  \theta_{k+1} = \Pi_\Theta(\theta_k + \kappa\,\Delta\theta_k),
  \label{eq:ch4:segment_update}
\end{equation}
where $J_k = \partial\vec{o}_k/\partial\theta$ is the segment-level
sensitivity Jacobian, $\lambda_J$ is a Tikhonov damping term, and
$\Pi_\Theta$ projects onto the feasible coefficient set
(e.g., $\alpha_j, \lambda_s, \lambda_h, \gamma \geq 0$).
This is a fast, local correction that does not require full
backpropagation through the episode.

\subsection{Episode Timescale: CVaR Gradient Descent}
\label{subsec:ch4:episode}

After each complete episode, the full loss \eqref{eq:ch4:L_total} is
differentiated through the rollout sensitivity recursion
\eqref{eq:ch4:dp_dtheta}--\eqref{eq:ch4:dq_dtheta} and the parameters
are updated:
\begin{equation}
  \theta \leftarrow \theta - \eta\,\nabla_\theta\,\mathcal{L}_{\mathrm{total}}.
  \label{eq:ch4:episode_update}
\end{equation}
Because $\mathcal{L}_{\mathrm{nav}}$ is a CVaR loss, the gradient
concentrates on the worst-performing episodes in the batch, focusing
learning effort on failure modes rather than average behavior.

\subsection{Curriculum Timescale: Passivity-Guided Advancement}
\label{subsec:ch4:curriculum}

The outermost loop governs when the learner transitions from
Phase~1 (geometry-only) to Phase~2 (material-aware), and which
start–goal pairs are sampled at each stage.

\paragraph{Passivity criterion.}
The discrete passivity surrogate is:
\begin{equation}
  \mathcal{L}_{\mathrm{pass}}
  = \frac{1}{T}\sum_t \mathrm{softplus}(\Delta H_t),
  \label{eq:ch4:L_pass}
\end{equation}
where $\Delta H_t = H_\theta(\PhaseState_{t+1}) - H_\theta(\PhaseState_t)$.
$\mathcal{L}_{\mathrm{pass}} \approx 0$ means the stored energy is
non-increasing (passive); large values indicate the learner is
injecting energy—a sign of instability.

\paragraph{Coverage criterion.}
A start–goal pair $(\vec{q}_0, \vec{g})$ is considered
\emph{conquered} in episode $e$ if:
\begin{equation}
  \hat{c}_e(\vec{q}_0,\vec{g})
  = \mathbf{1}\!\left[
      \mathcal{L}_{\mathrm{pass}} \leq \epsilon_{\mathrm{pass}}
      \;\wedge\;
      |\vec{q}_T - \vec{g}| \leq \epsilon_{\mathrm{goal}}
    \right].
  \label{eq:ch4:conquest}
\end{equation}
The overall coverage at episode $e$ is:
\begin{equation}
  \mathcal{C}(\mathcal{D}_e)
  = \frac{1}{|\mathcal{D}_e|}
    \sum_{(\vec{q}_0,\vec{g})\in\mathcal{D}_e}
    \hat{c}_e(\vec{q}_0,\vec{g}).
  \label{eq:ch4:coverage}
\end{equation}

\paragraph{Sampling and advancement.}
The curriculum re-weights the sampling distribution toward
unconquered pairs:
\begin{equation}
  \Pi_{\mathrm{curr}}^{e+1}(\vec{q}_0,\vec{g})
  \;\propto\; 1 - \hat{c}_e(\vec{q}_0,\vec{g}).
  \label{eq:ch4:sampling}
\end{equation}
Phase advancement from P1 $\to$ P2 occurs when:
\begin{equation}
  \mathcal{L}_{\mathrm{pass}} \leq \epsilon_{\mathrm{pass}}
  \quad\text{and}\quad
  \mathcal{C}(\mathcal{D}_e) \geq \rho_{\mathrm{cover}}.
  \label{eq:ch4:advance}
\end{equation}
This criterion answers the question: \emph{when has a region been
sufficiently conquered?}
The learner must be both passive (stable) and broadly successful
(covering $\rho_{\mathrm{cover}}$ of the sampled configurations)
before the curriculum advances.

% -----------------------------------------------------------------------
\section{Algorithm}
\label{sec:ch4:algorithm}
% -----------------------------------------------------------------------

Algorithm~\ref{algo:ch4:material} operationalizes the three-timescale
structure.
The outer loop iterates over curriculum episodes; within each episode,
an inner segment loop executes the dynamics and local updates; after
each episode the CVaR gradient step and passivity check run.

\begin{algorithm}[t]
\caption{Material-Aware Continual Hamiltonian Learning}
\label{algo:ch4:material}
\begin{algorithmic}[1]
\State \textbf{Input:} $(\vec{q}_0,\vec{g})$, geometry $Z^{(1)}$,
       oracle material $Z^{(2)}$, phase $\phi\in\{\mathrm{P1},\mathrm{P2}\}$
\State \textbf{Init:} $\theta \leftarrow \theta_{\mathrm{geom}}$ (P1) or
       warm-started from P1 (P2); $\mathcal{D} \leftarrow$ initial
       start–goal distribution; $e \leftarrow 0$
\While{$e < E_{\max}$ \textbf{and} phase not completed}
  \State Sample $(\vec{q}_0,\vec{g}) \sim \Pi_{\mathrm{curr}}^e$
  \State Initialize $\PhaseState_0 = (\vec{q}_0, \mathbf{0})$;
         compute $\{g_k\}$ via geometry-only A$^*$ scaffold
  \State $J_{\mathrm{ep}} \leftarrow 0$,\;
         $\mathcal{L}_{\mathrm{pass}} \leftarrow 0$
  \For{segment $k = 1, \ldots, K$}
    \State Query $f_\theta(\mathcal{E}_k, P_k) \to
           (\beta,\{\alpha_j\},\lambda_s,\lambda_h,\gamma)$ \hfill
           $\triangleright$ CoefEnergyNetMaterial
    \For{step $t = 1, \ldots, T_k$}
      \State Compute $F_{\mathrm{tot}}$ via \eqref{eq:ch4:F_tot}
      \State $\Mom_{t+1} \leftarrow \Mom_t + \tau\,F_{\mathrm{tot}}$;\;
             $\vec{q}_{t+1} \leftarrow \vec{q}_t + \tau\Mass^{-1}\Mom_{t+1}$
      \State Accumulate $J_{\mathrm{ep}} \mathrel{+}= j_t$;\;
             $\mathcal{L}_{\mathrm{pass}} \mathrel{+}= \mathrm{softplus}(\Delta H_t)/T$
      \State Update observer: $y_{t+1} = \Psi(y_t,\vec{q}_{t+1})$
    \EndFor
    \State \textbf{Segment update:} $\theta \leftarrow \Pi_\Theta(\theta
           + \kappa\,\Delta\theta_k)$ via \eqref{eq:ch4:segment_update}
  \EndFor
  \State \textbf{Episode update:}
         $\theta \leftarrow \theta - \eta\,\nabla_\theta\mathcal{L}_{\mathrm{total}}$
         via \eqref{eq:ch4:episode_update}
  \State Mark $(\vec{q}_0,\vec{g})$ conquered if \eqref{eq:ch4:conquest}
         holds
  \State Update $\Pi_{\mathrm{curr}}^{e+1}$ via \eqref{eq:ch4:sampling}
  \If{\eqref{eq:ch4:advance} holds and $\phi = \mathrm{P1}$}
    \State Advance to P2: unfreeze $(\lambda_s,\lambda_h)$
  \EndIf
  \State $e \leftarrow e + 1$
\EndWhile
\State \textbf{Return:} learned $\theta$, trajectory archive
\end{algorithmic}
\end{algorithm}

% -----------------------------------------------------------------------
\section{Instantiation: CoefEnergyNetMaterial}
\label{sec:ch4:implementation}
% -----------------------------------------------------------------------

\paragraph{Architecture.}
\textbf{CoefEnergyNetMaterial} extends the geometry backbone of
Chapter~\ref{chap:grlsnam} with a risk encoder and two material
coefficient heads:
\begin{itemize}
  \item \textbf{Geometry branch.}
        A Transformer operating on the local obstacle token set
        $\{(d_j, \vec{n}_j)\}$, producing geometry embeddings
        $\vec{h}_{\mathrm{geom}}$.
        Outputs: $(\beta, \{\alpha_j\}, \gamma)$.
  \item \textbf{Risk CNN encoder.}
        A small CNN operating on the $6$-channel risk patch
        $P_t(\vec{q}) \in \mathbb{R}^{6\times P\times P}$,
        producing a risk embedding $\vec{h}_{\mathrm{risk}}$.
  \item \textbf{Material heads.}
        Two linear heads from $[\vec{h}_{\mathrm{geom}}; \vec{h}_{\mathrm{risk}}]$
        producing $\lambda_s \geq 0$ and $\lambda_h \geq 0$
        (via softplus activations).
\end{itemize}

\paragraph{Barrier definitions.}
The IPC-style geometry barrier is
$b_{\mathrm{IPC}}(d) = -\log(d/d_{\mathrm{hat}})$ for $d < d_{\mathrm{hat}}$,
$0$ otherwise.
The hazard SDF barrier is
$b_{\mathrm{sp}}(\phi) = \mathrm{softplus}(-\phi/\sigma_\phi)$
with $\sigma_\phi = 1\,\text{m}$.

\paragraph{Two-phase training.}
\begin{itemize}
  \item \textbf{Phase~1.}
        Train with $\mathcal{L}_{\mathrm{P1}}$; material heads
        $(\lambda_s, \lambda_h)$ are frozen at $0$.
        The curriculum advances when $\mathcal{L}_{\mathrm{pass}}
        \leq \epsilon_{\mathrm{pass}}$ and
        $\mathcal{C} \geq \rho_{\mathrm{cover}}$.
  \item \textbf{Phase~2.}
        Initialize from Phase~1 weights; unfreeze
        $(\lambda_s, \lambda_h)$; train with $\mathcal{L}_{\mathrm{P2}}$.
        The geometry backbone is retained as a scaffold.
\end{itemize}

% \paragraph{Phase~3 (deferred).}
% Learning the material map $\tilde{r}$ from raw RGB-IR imagery—replacing
% oracle access with learned perception—is the natural next step but
% is beyond the current chapter's scope.

% -----------------------------------------------------------------------
\section{What the Richer Learner Achieves}
\label{sec:ch4:results}
% -----------------------------------------------------------------------

The material-aware learner is validated on the DFC2018 Houston aerial
dataset, where surface materials (water, highways, railways) create
latent hazards invisible to LiDAR geometry.

\paragraph{Risk reduction.}
Across 20 held-out test episodes, Stage~2 reduces cumulative soft-risk
exposure by \textbf{56.4\%} relative to the geometry-only Stage~1
model and by \textbf{17.4\%} relative to a material-blind Dijkstra
baseline, closing 27\% of the gap to the oracle upper bound.

\paragraph{Hazard elimination.}
Hard-hazard incursion count drops by 91\% from Stage~1 to Stage~2:
the hard-hazard barrier $F_{\mathrm{hard}}$ successfully enforces
forward invariance of the hazard-free set, as predicted by the
port-Hamiltonian analysis.

\paragraph{Geometry retained.}
Path length and SPL are within 3\% of Stage~1 on geometry-only test
cases, confirming that the material risk heads do not degrade the
geometry backbone when $\tilde{r} \approx 0$.

\paragraph{Stability.}
The passivity criterion $\mathcal{L}_{\mathrm{pass}} \leq 0.01$ is
satisfied at convergence, and the limit-cycle oscillations observed
in Stage~1 on narrow passages are eliminated—consistent with the
separate dissipative port providing a clean damping channel.

% -----------------------------------------------------------------------
\section{Chapter Summary}
\label{sec:ch4:summary}
% -----------------------------------------------------------------------

This chapter enlarged the geometry-only learner of
Chapter~\ref{chap:grlsnam} along four coupled axes:

\begin{enumerate}
  \item \textbf{Stored energy.}
        $H_\theta$ now carries soft-risk and hard-hazard barrier terms
        in addition to geometry; the dissipative port $R_\theta$ is
        separated explicitly (\S\ref{sec:ch4:energy}).
  \item \textbf{Coefficient space.}
        Two new coefficients $(\lambda_s, \lambda_h)$ are added, each
        with its own gradient pathway through the rollout sensitivity
        recursion (\S\ref{sec:ch4:gradients}).
  \item \textbf{Objective.}
        A CVaR tail-risk loss replaces the expected-cost surrogate,
        directly optimizing rare material-hazard encounters
        (\S\ref{sec:ch4:objective}).
  \item \textbf{Adaptation hierarchy.}
        Three nested loops: segment Gauss-Newton, episode CVaR descent,
        and passivity-guided curriculum provide stable continual
        learning across timescales (\S\ref{sec:ch4:multiscale}).
\end{enumerate}

These four enrichments are coupled: each new energy term requires its
own force channel, coefficient, objective term, and update law.
The chain from $H_i \to F_i \to \vartheta_i \to \mathcal{L}_i \to
\text{update law}_i$ is the structural principle that makes this
a coherent richer learner rather than a collection of independent
additions.

Chapter~\ref{chap:hierarchy} will show that this chain is not specific
to the geometry-to-material transition, but is an instance of a general
\emph{enrichment principle} that can extend the framework further.